\chapter{Nasal Polyps}

\section*{Definition}
Nasal polyps are benign, edematous, grape-like projections of inflamed sinonasal mucosa arising from the paranasal sinuses (most commonly the ethmoid sinuses) and protruding into the nasal cavity.

\section*{What Happens in Your Body}
Chronic type 2 inflammatory response with predominant eosinophilic infiltration, TH2 immune signaling proteins polarization (IL-4, IL-5, IL-13), and local IgE production leads to mucosal edema, stromal hyperplasia, and epithelial proliferation forming polypoid tissue.

\section*{Causes}
\begin{itemize}
  \item \textbf{Chronic rhinosinusitis with nasal polyps (CRSwNP):} Most common; eosinophilic inflammatory subtype
  \item \textbf{Aspirin-exacerbated respiratory disease (AERD):} Triad of nasal polyps, asthma, and NSAID hypersensitivity (Samter triad)
  \item \textbf{Allergic fungal rhinosinusitis:} Type I hypersensitivity to fungal antigens with eosinophilic mucin
  \item \textbf{Genetic conditions:} Cystic fibrosis (CFTR mutations), primary ciliary dyskinesia (Kartagener syndrome)
  \item \textbf{Churg-Strauss syndrome (eosinophilic granulomatosis with polyangiitis):} Systemic vasculitis with polyps
\end{itemize}

\section*{When to See a Doctor}
See your doctor if:
\begin{itemize}
  \item Bilateral nasal obstruction not responsive to topical steroid therapy
  \item Anosmia (complete loss of smell) with nasal congestion
  \item Facial pressure or headache
  \item Recurrent acute sinusitis despite antibiotics
  \item History of asthma, aspirin sensitivity, or cystic fibrosis with nasal symptoms
  \item Unilateral polyp (suspect inverted papilloma or neoplasm)
\end{itemize}

\section*{Related Symptoms}
\begin{itemize}
  \item Nasal obstruction (complete bilateral often)
  \item Anosmia or hyposmia
  \item Rhinorrhea (anterior and postnasal)
  \item Facial pressure or headache
  \item Sleep disturbance and snoring
\end{itemize}
\textbf{Important:} Nasal polyps in children should raise suspicion for cystic fibrosis, and unilateral polyps in any age group must be biopsied to exclude inverted papilloma or malignancy.

\section*{Self-Care / Home Management}
\begin{itemize}
  \item Consistent use of intranasal corticosteroids (fluticasone, mometasone, budesonide) daily
  \item Saline nasal irrigation to improve mucosal clearance and topical medication penetration
  \item Avoidance of NSAIDs if known AERD
  \item Aspirin desensitization (under specialist care) for AERD
  \item Smoking cessation (smoking worsens polyp formation and surgical outcomes)
\end{itemize}

\section*{How Your Doctor Will Evaluate You}
\begin{itemize}
  \item Nasal endoscopy to visualize polyp location, extent, and bilaterality
  \item CT sinuses without contrast for mapping opacity extent and bony erosion
  \item MRI sinus if orbital or intracranial extension suspected
  \item Allergy evaluation (skin prick test, serum IgE) to identify triggers
  \item Sweat chloride testing or CFTR genetic testing in children and selected adults with polyps
  \item Biopsy for any unilateral, friable, or bleeding polyp
\end{itemize}

\section*{Treatment Options}
\begin{itemize}
  \item \textbf{First-line medical:} Intranasal corticosteroid sprays (fluticasone, mometasone) or drops (budesonide)
  \item \textbf{Second-line:} Short-course oral corticosteroids (prednisone 40 mg daily taper) for symptom flare
  \item \textbf{Biologic therapy:} Dupilumab (anti-IL-4R), omalizumab (anti-IgE), or mepolizumab (anti-IL-5) for refractory CRSwNP
  \item \textbf{Functional endoscopic sinus surgery (FESS):} Polypectomy and sinus ostia widening
  \item \textbf{Post-operative:} Saline irrigations and topical steroids; serial endoscopic debridement
  \item \textbf{Aspirin desensitization:} Maintenance aspirin therapy after graded challenge in AERD people
  \item \textbf{Management of comorbidities:} Asthma optimization, allergy immunotherapy
\end{itemize}

\section*{References}
\begin{thebibliography}{99}
\bibitem{nasal_polyps:polyp1} Fokkens WJ, Lund VJ, Hopkins C, et al. "European Position Paper on Rhinosinusitis and Nasal Polyps 2020." Rhinology. 2020;58(Suppl S29):1-464.
\bibitem{nasal_polyps:polyp2} Bachert C, Bhatt DL, Desrosiers M, et al. "Efficacy and safety of dupilumab in patients with severe chronic rhinosinusitis with nasal polyps (LIBERTY NP SINUS-24 and SINUS-52)." Lancet. 2019;394(10209):1638-1650.
\bibitem{nasal_polyps:polyp3} Orlandi RR, Kingdom TT, Hwang PH, et al. "International consensus statement on allergy and rhinology: rhinosinusitis." Int Forum Allergy Rhinol. 2016;6(Suppl 1):S22-S209.
\bibitem{nasal_polyps:polyp4} Stevens WW, Schleimer RP, Kern RC. "Chronic rhinosinusitis with nasal polyps." J Allergy Clin Immunol Pract. 2016;4(4):565-572.
\bibitem{nasal_polyps:polyp5} Larsen PL, Tos M, Baer S. "Nasal polyps." In: Cummings Otolaryngology. 7th ed. Elsevier; 2021.
\bibitem{nasal_polyps:polyp6} Laidlaw TM, Mullol J, Woessner KM, et al. "Chronic rhinosinusitis with nasal polyps and aspirin-exacerbated respiratory disease." J Allergy Clin Immunol Pract. 2022;10(1):63-73.
\end{thebibliography}
